 % TODO @Final: Zu lange Textzeilen beheben.
 % TODO @Final: Fußnoten prüfen!
 % TODO @Final: Pagebreaks prüfen!

\chapter{Einleitung}
\label{chap:einleitung}

\section{Hintergrund und Motivation}
\label{sec:motivation}
 
\emph{\enquote{Nichts ist so beständig wie der Wandel.}} -- Dieses Zitat des griechischen Philosophen Heraklit von Ephesus (535-475 v. Chr.) ist in der heutigen industriellen Praxis aktueller denn je. Produzierende Unternehmen der Prozessindustrie sehen sich mit einem immensen Wandel konfrontiert \cite{Valluru.2024}: Die Globalisierung intensiviert den internationalen Wettbewerb \cite{Scheuermann.2025}, während technologischer Fortschritt und Digitalisierung die Produktentwicklung beschleunigen \cite{Deloitte.2023}. Kunden erwarten zunehmend maßgeschneiderte Lösungen \cite{Gro.2022} und nachhaltige Produkte \cite{RuthStrau.2024,Deloitte.2023}. Gleichzeitig ändern sich regulatorische Rahmenbedingungen ständig \cite{Uniquelifescience.2025}, und globale Lieferketten erweisen sich als fragil gegenüber Naturkatastrophen, politischen Konflikten und Pandemien \cite{Bossut.2025}.

Diese Entwicklungen resultieren in volatilen Marktbedingungen, die Unternehmen dazu zwingen, Produktlebenszyklen sowie die Time-to-Market ihrer Produkte zu verkürzen, um wettbewerbsfähig zu bleiben \cite{Lier.2015.DE}. Dies wird begleitet von einer immer stärkeren Produktindividualisierung, erhöhter Produktvielfalt und schwankenden Produktionsmengen. Während ab dem Jahr 1913 im Zuge der Fließbandfertigung ein einziges Produkt in großer Stückzahl schnell hergestellt werden sollte, besteht heute in einigen Teilen der Produktion das Ziel der \emph{Mass Customization}, d.~h. der schnellen Fertigung individueller Produkte in sehr geringen Losgrößen \cite{Kirchgeorg.2018}.

Diese Trends verlangen Produktionsanlagen hinsichtlich ihrer hardware- und softwareseitigen Gestaltung ein hohes Maß an Anpassbarkeit ab -- im Sinne von Flexibilität und Wandelbarkeit. Der Begriff \emph{Flexibilität} beschreibt die Anpassbarkeit eines Produktionssystems an äußere oder innere Veränderungen (z.~B. neue Produkte) \emph{ohne} strukturelle Änderungen des Systems \cite{VDI.5201-1}. \emph{Wandelbarkeit} hingegen beschreibt die Anpassbarkeit eines Produktionssystems an äußere oder innere Veränderungen (z.~B. neue Produkte) \emph{mit} strukturellen Änderungen des Systems \cite{VDI.5201-1}.

Ein wesentlicher Befähiger, um die erforderliche Flexibilität und Wandelbarkeit der Produktionssysteme zu gewährleisten, ist die \emph{Modularisierung} dieser Systeme \cite{Lier.2015.DE,GroePuppendahl.2016,Seibold.2020}. Dies umfasst zum einen die Aufteilung der Produk\-tions\-systeme in einzelne weitestgehend voneinander unabhängige modulare Einheiten, aber auch deren Integration in ein modulares Gesamtsystem im Sinne eines individuellen Anwendungsfalls. Dies ermöglicht Wandelbarkeit, da für verschiedene Anwendungsfälle modulare Einheiten in unterschiedlicher Zusammenstellung genutzt werden können. Flexibilität entsteht dadurch, dass die modularen Einheiten selbst durch Parametrierung an verschiedene Anwendungsfälle angepasst werden können.

Um die Vorteile der physischen Modularisierung von Produktionssystemen wirtschaftlich nutzen zu können, ist es notwendig, dass auch deren Automatisierungstechnik modular und damit flexibel und wandelbar gestaltet ist. Die funktionale Aufteilung und vor allem die schnelle und einfache Integration modularer Einheiten stellen aus automatisierungstechnischer Sicht Herausforderungen dar \cite{Urbas.2012}. Vor diesem Hintergrund wurde im Umfeld der Prozessindustrie in verschiedenen Forschungsprojekten \cite{DECHEMA.,CORDIS.2021} und Standardisierungsarbeiten \cite{PNO.2025,VDIVDEGMA.DE,VDI.2020} das sogenannte \emph{Module~Type~Pa\-cka\-ge-Konzept (MTP-Konzept)} für die Automatisierung modularer Produktionsanlagen entwickelt. Dieses entstand beim \gls{vdi} als Richtlinienreihe \cite{VDIVDEGMA.DE} und wird mittlerweile durch die \gls{pno} als Richtlinienreihe \cite{PNO.2025} spezifiziert. 

Intelligente Module, sogenannte \emph{Process Equipment Assemblies (PEAs)}, werden hardware- und softwareseitig so vorbereitet und getestet, dass sie verfahrenstechnische Funktionen eigenständig umsetzen können. Aus den PEAs kann anschließend eine modulare Produktionsanlage für einen konkreten Anwendungsfall zusammengestellt werden. Hierfür muss die Hardware der PEAs manuell aufgebaut und physisch verbunden werden. Die automatisierten Funktionen der PEAs können mittels MTP-Konzept automatisiert zu einer übergreifenden Steuerung der modularen Gesamtanlage zusammengefügt werden. Die Richtlinienreihe \cite{PNO.2025} spezifiziert dafür standardisierte Automatisierungsschnittstellen zur Interaktion mit PEAs. Diese Schnittstellen werden in einer gepackten AutomationML-Datei, dem \emph{Module Type Package (MTP)}, beschrieben und mit Semantik hinterlegt. Die MTPs der PEAs mit den Schnittstelleninformationen können anschließend in ein übergeordnetes Automatisierungssystem, den \emph{Process Orchestration Layer (POL)}, eingelesen werden, der PEA-übergreifende Verwaltungs-, Koordinations- und Überwachungsaufgaben übernimmt. Damit ermöglicht das MTP-Konzept eine automatisierte und damit schnelle und einfache Integration intelligenter modularer Einheiten (PEAs) verschiedener Hersteller in ein modulares Produktionssystem.

Unter Anwendung des MTP-Konzepts können modulare Prozessanlagen schnell, einfach und damit wirtschaftlich auf- und umgebaut werden, was das Ziel der Wandelbarkeit unterstützt \cite{Holm.2016,NAMURProcessNetVDMAundZVEI.2021.DE}. Zudem ermöglichen MTP-Konzepte, wie z.~B. die Nutzung von Diensteparametern und Prozeduren, die Flexibilisierung der PEA-Funktion selbst. Auf diese Weise bildet das MTP-Konzept eine wesentliche Grundlage, um die bestehenden Flexibilitäts- und Wandelbarkeitsziele zu erreichen. 

Die durch das MTP-Konzept ermöglichte Anpassbarkeit von Produktionssystemen wirkt sich allerdings nicht nur auf die Produktionsanlagen selbst, sondern auch auf die produktionsnahe Versorgungs- und Verpackungslogistik aus. U.~a. die Autoren von \cite{Kessler.2015,Lier.2015.DE} zeigen, dass zum effizienten Betrieb modularer Produktionssysteme auch eine ebenso modulare und damit flexible und wandelbare produktionsnahe Logistik notwendig ist.

\begin{graybox}
	\textbf{Fazit:}	Die zunehmende Marktvolatilität und Produktindividualisierung sowie verkürzte Innovationszyklen erfordern flexible und wandelbare Produktions- und Logistiksysteme. Die Modularisierung der Hardware und Automatisierungssoftware dieser Systeme ist ein wesentlicher Befähiger für die erforderliche Flexibilität und Wandelbarkeit.
\end{graybox}

\section{Problemstellung}
\label{sec:problem}

Aus physischer Sicht weisen produktionsnahe Logistiksysteme bereits heute meist eine modulare Struktur auf und bestehen aus mehreren, hierarchisch angeordneten System\-ebenen. Die Basis bilden hochintegrierte Logistikmaschinen, sogenannte \gls{pu}. Jede PU ist hardware- und softwareseitig durch einen PU-Hersteller vorbereitet, um eine bestimmte produktionslogistische Funktion, wie z.~B. Fördern, Abfüllen oder Palettieren, zu erfüllen. Teils besitzen PUs eigene speicherprogrammierbare Steuerungen (engl. \gls{plc}), teils ist aus Kostengründen die Steuerungslogik mehrerer PUs auf einer PLC zusammengefasst. Obwohl PUs ihre Kernfunktionalität sehr eigenständig ausführen können, besitzen sie dennoch häufig Abhängigkeiten voneinander, z.~B. findet ein Austausch von Freigabesignalen statt. Zum Teil kommen Liniensteuerungen (engl. Line Controller) zum Einsatz, die mehrere PUs und deren Abhängigkeiten nach einer festprogrammierten Logik koordinieren. Daneben können \gls{mfr} die Koordination notwendiger Transporte und die Auswahl zu nutzender Transportressourcen übernehmen. \cite{NAMURAK4.19.2023,Lenz.2017}

PUs, Line Controller und MFRs können zudem in ein übergeordnetes Automatisierungssystem, meist ein \gls{mes}, ein \gls{les}, ein \gls{wms} oder ein \gls{pls}, eingebunden werden. Diese Systeme übernehmen logistiksystemweite Verwaltungs-, Koordinations- und Überwachungsaufgaben und dienen als Verbindungsschicht zu den betriebswirtschaftlich geprägten übergeordneten \gls{erp}-Systemen. \cite{Kilger.2006,Munnemann.2008,Lenz.2017}

In der Vergangenheit wurden für die Integration von PUs, Line~Controllern und MFR in die übergeordneten Automatisierungssysteme vielfältige firmen- und projekt\-spezifische Integrationslösungen entwickelt. Bei der Integration gilt es, die zu integrierenden Einheiten an die Gegebenheiten des Systems vor Ort anzupassen und verschiedenste Informationsflüsse zwischen den Einheiten untereinander und zum übergeordneten Automatisierungssystem umzusetzen. So werden im Falle der PU-Integration dem übergeordneten Automatisierungssystem beispielsweise Maschinen- und Prozessdaten zur Auswertung bereitgestellt. Umgekehrt übernimmt dieses die Parametrierung der PUs und die Koordination der technischen Abläufe im System. Dies erfordert umfangreiche Abstimmungen der zu übertragenden Informationen zwischen den Herstellern der verschiedenen PUs, dem Hersteller des übergeordneten Automatisierungssystems und dem Anwender. Die erforderlichen Abstimmungen basieren meist auf firmen- oder projektspezifischen Listen. Darin werden alle zu übertragenden Informationen, deren Datentypen und ggf. für die Kommunikation notwendige Angaben, wie z.~B. der Telegrammaufbau, festgehalten. Häufig kommen in heutigen Logistiksystemen digitale Kommunikationstechnologien, wie z.~B. PROFIBUS oder PROFINET, in einigen Fällen aber auch analoge Signale zum Einsatz. \cite{Gunther.2003,Kilger.2006,Lenz.2017,NAMURAK4.19.2023}

Proprietäre Schnittstellen und Dokumentation sowie unterschiedliche Informationsübertragung verursachen erhebliche Aufwände hinsichtlich firmenübergreifender Abstimmungen, Implementierungen und Tests beim Auf- und Umbau produktionsnaher Logistiksysteme. Folgen sind lange Planungszeiten, eine schlechte Wiederverwendbarkeit von Logistikmaschinen und damit verbundene hohe Kosten. Aufgrund der hohen Abstimmungsaufwände werden zudem nur die notwendigsten Informationen zwischen den PUs untereinander und zu übergeordneten Automatisierungssystemen ausgetauscht. Dieser Datenaustausch beschränkt sich zudem häufig auf einfach strukturierte Datenobjekte, z.~B. Messwerte oder Steuerungsgrößen. Weiterführendes Potenzial der in den PUs vorhandenen Informationen, z.~B. zur genaueren Prozessüberwachung, bleibt ungenutzt. \cite{NAMURAK4.19.2023,Lenz.2017}

Aufgrund dieser Gegebenheiten können trotz des hardwareseitig modularen Aufbaus produktionsnaher Logistiksysteme die durch die modulare Produktion geforderte Flexibilität und Wandelbarkeit nicht erreicht werden. Zudem müssen die notwendigen Abstimmungen durch Experten durchgeführt werden, die das Verpackungssystem und den umzusetzenden Prozess sehr gut kennen. Solche Experten werden in Zukunft immer weniger zur Verfügung stehen \cite{Pieringer.2025}. 

Diese Problemstellungen werden auch im Positionspapier \cite{NAMURAK4.19.2023} des \gls{namur}-Arbeitskreises 4.19 \enquote{Produktionsnahe Logistik} adressiert. Es wird eine automatisierungstechnische Integrationslösung produktionslogistischer Module gefordert, die sowohl alle technischen Anforderungen (z.~B. gemäß \cite{NAMURAK4.19.2020}) erfüllt als auch den zeit- und kostenintensiven Abstimmungs- und Anpassungsaufwand signifikant reduziert. Dazu werden Standards benötigt, die für alle zu integrierenden Modulinformationen genau eine Lösung interpretationsfrei spezifizieren und Abstimmungen obsolet machen. Dadurch soll eine standardisierte und automatisierte Integration produktionslogistischer Module ermöglicht werden. 

Derartige Integrationsstandards versprechen erhebliche Kosteneinsparungen im Bereich der produktionsnahen Logistik. Gemäß \cite{NAMURAK4.19.2023} betragen die Integrationskosten für Logistikmodule in produzierenden Unternehmen heute ca. 156.000~\euro~pro Unternehmen und Jahr unter der Annahme von durchschnittlich 30 Modulintegrationen. Zukünftig wird im Kontext der modularen Produktion und der Digitalisierung ein deutlicher Anstieg dieser Kosten, u.~a. durch eine höhere Anzahl von Modulintegrationen, erwartet. Darüber hinaus steigen auch die Kosten für Implementierung, Tests und Dokumentation aufseiten der Modulhersteller, was zu höheren Anschaffungskosten für Logistikmodule führt. Automatisierte Integrationsprozesse und standardisierte Datenstrukturen bieten einen Weg zur Eindämmung dieser Kosten, sowohl für Modulhersteller als auch für produzierende Unternehmen. \cite{NAMURAK4.19.2023}

\begin{graybox}
	\textbf{Fazit:}	Obwohl produktionsnahe Logistiksysteme hardwareseitig modular aufgebaut sind, verhindern proprietäre Schnittstellen, aufwändige manuelle Abstimmungs- und Engineeringprozesse sowie fehlende Integrationsstandards die Erreichung der Flexibilitäts- und Wandelbarkeitsziele.
\end{graybox}

\section{Zielstellung und Forschungsfragen}
\label{sec:ziel}

Vor dem Hintergrund dieser Problemstellung erarbeitete der NAMUR-Arbeitskreis~4.19 die NAMUR-Empfehlung \cite{NAMURAK4.19.2020}, welche Anforderungen an die Automatisierung modularer produktionsnaher Logistiksysteme der Prozessindustrie beschreibt. Zudem zeigten die Autoren von \cite{NAMURAK4.19.2020} und die Autoren von \cite{Cordes.2020.DE} anhand einer Simulationsstudie, dass das MTP-Konzept grundsätzlich für die Automatisierung modularer produktionsnaher Logistiksysteme geeignet ist.

Allerdings erfüllt das bisher auf die modulare Produktion der Prozessindustrie ausgerichtete MTP-Konzept nicht alle Anforderungen der produktionsnahen Logistik (z.~B. gemäß \cite{NAMURAK4.19.2020}). Beispielsweise ist eine direkte Modul-zu-Modul-Kommunikation, die zwischen Logistikmaschinen notwendig ist, heute in \cite{PNO.2025} nicht vorgesehen. Ebenso fehlt die Möglichkeit zur Integration fahrerloser Transportsysteme, die für flexible Logistiksysteme unabdingbar sind. 

An dieser Stelle setzt die vorliegende Dissertation an und beschreibt ein Konzept zur Automatisierung modularer produktionsnaher Logistiksysteme. Ziel ist, durch eine einheitliche, automatisierte Integration von Logistikmodulen die Ab\-stim\-mungs- und Engineeringaufwände im Zuge des Auf- und Umbaus modularer produktionsnaher Logistiksysteme signifikant zu reduzieren. Dafür werden die in \cite{PNO.2025} standardisierten sowie im MTP-Umfeld erforschten Konzepte herangezogen und derart weiterentwickelt, dass sie die Anforderungen modularer produktionsnaher Logistiksysteme (z.~B. gemäß \cite{NAMURAK4.19.2020}) erfüllen. 

Die Inhalte dieser Dissertation wurden maßgeblich im Rahmen des \gls{enpro} 2.0-Projekts \gls{moprolog} erarbeitet, welches zum Ziel hatte, Anlagenkonzepte für modulare produktionsnahe Logistiksysteme zu entwickeln. Insbesondere wurde die Automatisierung von Logistikmodulen sowie deren Integration in ein modulares Gesamtsystem im Kontext der \textbf{Verpackungslogistik} betrachtet.

Aus der beschriebenen Zielstellung und dem Fokus des MoProLog-Projekts resultieren zwei Forschungsfragen. Diese sind im Folgenden gemäß der Design Science Research-Methode nach Wieringa \cite{Wieringa.2014} formuliert und enthalten folgende Bestandteile: Artefakt (\textcolor{BurntOrange}{Orange}), Problemkontext (\textcolor{NavyBlue}{Blau}), Anforderungen (\textcolor{OliveGreen}{Grün}) und Stakeholderziele (\textcolor{Maroon}{Rot}). Details zu Design Science Research folgen in Abschnitt~\ref{sec:methodik}.

\begin{graybox}
	\textbf{Forschungsfrage 1:} Wie sieht \textcolor{BurntOrange}{ein Automatisierungskonzept} für \textcolor{NavyBlue}{modulare produktionsnahe Logistiksysteme mit Fokus auf Verpackungssysteme} aus, das \textcolor{OliveGreen}{bestehende sowie in der Standardisierung befindliche Mechanismen, Schnittstellen und Modelle des Module~Type~Package-Konzepts nutzt}, um \textcolor{Maroon}{flexible und wandelbare Logistikprozesse umzusetzen}? 
\end{graybox}

\begin{graybox}
	\textbf{Forschungsfrage 2:} Welche \textcolor{BurntOrange}{neuen Mechanismen, Schnittstellen und Modelle} sind für \textcolor{NavyBlue}{modulare produktionsnahe Logistiksysteme mit Fokus auf Verpackungssysteme} notwendig, die \textcolor{OliveGreen}{auf dem bestehenden Module~Type~Package-Konzept aufsetzen} und \textcolor{Maroon}{eine Realisierung flexibler und wandelbarer Logistikprozesse ermöglichen}?  
\end{graybox}

\section{Forschungsmethodik}
\label{sec:methodik}

Diese Dissertation folgt einer Forschungsmethodik des \gls{dsr}. In diesem Abschnitt werden die Grundlagen von DSR (Abschnitt~\ref{sec:DSR}) und die verwendete Methode nach Dresch et al. (Abschnitt~\ref{Dresch}) vorgestellt.

\subsection{Design Science Research}
\label{sec:DSR}

Gemäß \cite{Wieringa.2014} ordnet sich DSR zwischen Grundlagenforschung und Ein\-zel\-fall-bezogener Forschung ein. Die Grundlagenforschung (z.~B. Physik, Chemie und Biologie) untersucht Naturgesetze, die überall Gültigkeit haben. Sie basiert auf universellen Verallgemeinerungen und idealisierten Annahmen, die umfassende Forschungsprobleme handhabbar machen, aber die reale Welt nur zu einem gewissen Grad abbilden. Die Einzelfall-bezogene Forschung (z.~B. Engineering, Beratung und Gesundheitswesen) untersucht konkrete Probleme der realen Welt und strebt keine Generalisierung über den Einzelfall hinaus an. Die Experten bauen Wissen auf, das verallgemeinert werden kann, aber nicht universell gültig ist. \cite{Wieringa.2014}

% Problemklassen
Design Science Research liegt im Bereich zwischen diesen beiden Extremen. Es beschäftigt sich mit Lösungen, die zwar über mehrere (möglichst viele) Einzelfälle hinweg generalisierbar sind, aber dennoch keine universelle Gültigkeit haben. Diese Lösungen sind für eine bestimmte Klasse von Problemen einsetzbar, die in ihren Charakteristika ähnlich sind. Zudem werden keine unrealistischen Abstraktionen und Idealisierungen vorgenommen, sondern realistische und begründete Annahmen über die Problemklasse getroffen. \cite{Wieringa.2014}

% Artefakte und Artefaktarten
Das Ziel von DSR ist, Artefakte zu entwerfen und zu bewerten, die eine Problemstellung in einem bestimmten Problemkontext bzw. einer bestimmten Einsatzumgebung lösen und dabei die Anforderungen relevanter Stakeholder erfüllen \cite{Wieringa.2014}. Artefakte können gemäß \cite{Dresch.2015} den in Tabelle~\ref{tab:Artefaktarten} dargestellten Arten zugeordnet werden, an denen sich auch die Artefakte dieser Dissertation orientieren (siehe Kapitel~\ref{chap:KonzeptUebersicht}). Während die Grundlagenforschung darauf abzielt, die Wahrheit zu finden, zielt DSR auf die Nützlichkeit der entwickelten Artefakte im betrachteten Problemkontext ab \cite{deSordi.2021}. Die Artefakte interagieren mit dem Problemkontext, der Problemkontext selbst kann aber nicht verändert werden \cite{Wieringa.2014}.

% Artefaktarten
\begin{footnotesize}
	\begin{longtable}{|
		>{\columncolor[HTML]{E0E0E0}}l 
		>{\columncolor[HTML]{FFFFFF}}l |}
		\caption{Artefaktarten in Design Science Research gemäß \cite{Dresch.2015}}
		\label{tab:Artefaktarten}\\
		\hline

		\multicolumn{1}{|L{2.5cm}|}{\cellcolor[HTML]{E0E0E0} \textbf{Bezeichnung}}        	& 
		\multicolumn{1}{L{12.13cm}|}{\cellcolor[HTML]{E0E0E0}\textbf{Beschreibung}} 
		\\ \hline
	
		\multicolumn{1}{|L{2.5cm}|}{\cellcolor[HTML]{FFFFFF} Konstrukt (auch: Konzept) }        	& 
		\multicolumn{1}{L{12.13cm}|}{\cellcolor[HTML]{FFFFFF} Diese Artefakte beschreiben Probleme innerhalb einer Domäne und spezifizieren Lösungen für diese. Sie können als das Vokabular einer Domäne verstanden werden.} 
		\\ \hline

		\multicolumn{1}{|L{2.5cm}|}{\cellcolor[HTML]{FFFFFF} Modell }        	& 
		\multicolumn{1}{L{12.13cm}|}{\cellcolor[HTML]{FFFFFF} Diese Artefakte beschreiben (Teil-) Konstrukte und vor allem deren Beziehungen zueinander. Der Fokus liegt in DSR auf dem Nutzen der Modelle, nicht auf einer möglichst hohen Genauigkeit.} 
		\\ \hline

		\multicolumn{1}{|L{2.5cm}|}{\cellcolor[HTML]{FFFFFF} Methode }        	& 
		\multicolumn{1}{L{12.13cm}|}{\cellcolor[HTML]{FFFFFF} Diese Artefakte beschreiben Schritte, die zur Ausführung bestimmter Aufgaben erforderlich sind. Methoden können an Modelle geknüpft sein und diese als Input verwenden. } 
		\\ \hline

		\multicolumn{1}{|L{2.5cm}|}{\cellcolor[HTML]{FFFFFF} Instanziierung }        	& 
		\multicolumn{1}{L{12.13cm}|}{\cellcolor[HTML]{FFFFFF} Diese Artefakte stellen die Realisierung und Operationalisierung anderer Artefakte (Konstrukte, Modelle, Methoden) in ihrer realen Umgebung dar. So wird die Wirksamkeit dieser Artefakte und deren Anwendungsweise in der Praxis gezeigt. } 
		\\ \hline

		\multicolumn{1}{|L{2.5cm}|}{\cellcolor[HTML]{FFFFFF} Gestaltungs-vorschlag }        	& 
		\multicolumn{1}{L{12.13cm}|}{\cellcolor[HTML]{FFFFFF} Diese Artefakte entsprechen einer generischen Vorlage, die verwendet werden kann, um Lösungen für eine bestimmte Problemklasse zu entwickeln. Sie beziehen sich somit auf den Forschungsprozess selbst. } 
		\\ \hline

		
	\end{longtable}
\end{footnotesize}

Das Ziel von DSR erfordert ein zweigeteiltes Vorgehen, welches in \cite{Wieringa.2014} durch einen \emph{Design Cycle} und einen \emph{Empirical Cycle} beschrieben wird. Der \emph{Design Cycle} befasst sich mit dem kreativen Schaffen notwendiger Artefakte. Im \emph{Empirical Cycle} werden Wissensfragen über die entwickelten Artefakte beantwortet, vorrangig mit dem Ziel, deren Performance im betrachteten Problemkontext zu evaluieren. 

In dieser Arbeit bietet sich der Einsatz von DSR-Methoden an, da systematisch nützliche Automatisierungskonzepte (= Artefakte) für den Einsatzbereich der Verpackungslogistik (= Problemkontext) geschaffen werden sollen. Dabei ist es wichtig zu beachten, dass DSR die Forschung unterstützen und niemals einschränken soll. Es sind passende DSR-Methoden für die vorliegende Problemstellung auszuwählen, zu interpretieren und nach Bedarf sinnvoll zu kombinieren.

\subsection{Methode nach Dresch et al.}
\label{Dresch}

Abbildung~\ref{DreschFig} zeigt die in dieser Dissertation vorrangig verwendete DSR-Methode nach Dresch et al. \cite{Dresch.2015}. Die Methode besteht aus 12 Schritten, die nacheinander, bei Bedarf auch iterativ, durchlaufen werden. In \emph{Identification of the problem} werden das zu lösende Problem sowie dessen Relevanz beschrieben und Forschungsfragen abgeleitet. Anschließend folgt der Schritt \emph{Awareness of the problem}, bei dem das Problem und der Problemkontext mit dem Ziel untersucht werden, alle relevanten Aspekte des Problems sowie dessen Grenzen zu ermitteln. In diesem Zusammenhang wird maßgeblich auch bestehende Literatur einbezogen (\emph{Systematic literature review}). Es folgt \emph{Identification of the artifacts and configuration of the classes of problems}. In diesem Schritt werden die Problemklassen und bereits bestehende Artefakte untersucht, die im Zusammenhang mit dem zu lösenden Problem stehen. Daraus werden generische Artefakte abgeleitet, die zur Lösung des Problems notwendig sind, und Performance-Kriterien für diese festgelegt. 

\begin{figure}[htbp]
	\centerline{\includegraphics[width=1.0\textwidth]{Inhalt/Abbildungen/01_Einleitung/Dresch.PNG}}
	\caption{Design Science Research-Methode nach Dresch et al. \cite{Dresch.2015}}
	\label{DreschFig}
\end{figure}

Auf Basis dieses Wissens werden in \emph{Proposition of artifacts to solve a specific problem} konkrete Artefakte zur Lösung des Problems im vorliegenden Problemkontext identifiziert. Dabei können auch mehrere alternative Artefakte vorgeschlagen werden. Anschließend werden aus diesen Artefaktvorschlägen die vielversprechendsten ausgewählt. Diese werden in \emph{Design of the selected artifacts} im Detail entworfen und in ihrer Funktion sowie ihren Möglichkeiten zur Umsetzung beschrieben. In \emph{Development of the artifact} findet die eigentliche Entwicklung der Artefakte in ihrer funktionalen Form (z.~B. Softwareimplementierung) statt. Im Schritt \emph{Evaluation of the artifact} erfolgt auf Basis von Beobachtungen und Messungen eine Bewertung der Artefakte im möglichst realitätsnahen Einsatzbereich. Bei der Bewertung ist auf die Erkenntnisse aus dem Schritt \emph{Awareness of the problem} zurückzugreifen und Wissensfragen in Bezug auf die Performance der Artefakte zu beantworten. 

Nach der Evaluation werden in \emph{Clarification of learning achieved} Aspekte beschrieben, die positiv und negativ auf den Forschungsprozess eingewirkt haben. Daran können sich andere Forscher in ähnlichen Bereichen orientieren. Anschließend werden die Ergebnisse der Forschung, deren Limitierungen und die während des Forschungsprozesses getroffenen Entscheidungen im Schritt \emph{Conclusion} zusammengefasst. In \emph{Generalization for a class of problems} wird gezeigt, dass die erarbeiteten Artefakte nicht nur für ein ganz bestimmtes Problem, sondern für eine Klasse von Problemen einsetzbar sind. Schließlich sind die Ergebnisse der Forschung im letzten Schritt \emph{Communication of the results} zu veröffentlichen und somit allen Interessierten zugänglich zu machen.

Im Folgenden wird ein Überblick darüber gegeben, welche Kapitel dieser Dissertation welche Teile der zuvor beschriebenen DSR-Methode behandeln.

\section{Aufbau der Arbeit}
\label{chap:aufbau}

Die Abschnitte~\ref{sec:motivation} bis \ref{sec:ziel} stellten bereits im Sinne des Schritts \emph{Identification of the problem} die Automatisierung modularer Verpackungssysteme als Problemstellung dieser Dissertation sowie deren Relevanz vor. Nach einer Einführung wesentlicher Grundlagen zum Verständnis dieser Arbeit in Kapitel~\ref{chap:Grundlagen} wird in Kapitel~\ref{chap:problemkontext} die detaillierte Analyse des Problemkontexts (\emph{Awareness of the problem}) behandelt. Daraus resultieren eine Struktur und Arbeitsweise modularer Logistiksysteme sowie Anforderungen an deren Automatisierung. Kapitel~\ref{chap:sota} zeigt den Stand der Erkenntnisse als Ergebnis des \emph{Systematic literature review} und führt in relevante Problemklassen ein. Aus einer Bewertung des Stands der Erkenntnisse wird zudem die Forschungslücke herausgearbeitet, die in dieser Dissertation adressiert wird.

\pagebreak

Kapitel~\ref{chap:KonzeptUebersicht} gibt einen Überblick über die zu entwickelnden Artefakte sowie deren Artefaktarten und ordnet diese Artefakte in die identifizierten Problemklassen ein (\emph{Identification of the artifacts and configuration of the classes of problems}). Die anschließenden Kapitel behandeln Konzepte zur Automatisierung und Integration von Logistikmodulen (Kapitel~\ref{chap:Art1LEA}), von Logistiklinien (Kapitel~\ref{chap:Art2LL}) und von flexiblen Transportsystemen in Logistikbereichen (Kapitel~\ref{chap:Art3LA}). Diese Konzeptkapitel stellen die Ergebnisse der Schritte \emph{Proposition of artifacts to solve a specific problem} und \emph{Design of the selected artifacts} dar. 

Kapitel~\ref{chap:evaluation} behandelt die praktische Umsetzung der entwickelten Artefakte anhand dreier Evaluationsbeispiele aus der Industrie, die jeweils in einem (Soft\-ware-) Demonstrator umgesetzt wurden (\emph{Development of the artifact}). Auf dieser Grundlage werden die entwickelten Artefakte anhand der eingangs identifizierten Anforderungen bewertet und der Nutzen der Artefakte zur Lösung der Problemstellung beschrieben (\emph{Evaluation of the artifact}). Zudem werden anhand der drei Evaluationsbeispiele Potenziale und Limitierungen der entwickelten Artefakte gezeigt (\emph{Clarification of learning achieved}).

Anschließend fasst Kapitel~\ref{chap:fazit} die erlangten Erkenntnisse zusammen und beantwortet die Forschungsfragen (\emph{Conclusion}). Kapitel~\ref{chap:ausblick} zeigt schließlich wei\-ter\-füh\-ren\-de Forschungs-, Entwicklungs- und Standardisierungspotenziale. Dabei wird u.~a. auf die Anwendbarkeit der Artefakte in der modularen Produktion der Fer\-ti\-gungs- und Prozessindustrie eingegangen (\emph{Generalization for a class of problems}). 

In Bezug auf den Schritt \emph{Communication of the results} wurden die erarbeiteten Ergebnisse in Konferenzbeiträgen, Fachmagazinen und einem Fachbuch veröffentlicht (Literaturverzeichnis, Abschnitt \enquote{Publikationen der Autorin}). Zudem wurden die Forschungsergebnisse auf Industriemessen, wie der Hannover Messe und der ACHEMA, dem interessierten Fachpublikum vorgestellt sowie in die Standardisierungsarbeiten zum MTP-Konzept beim VDI und der PNO eingebracht.